\PassOptionsToPackage{unicode}{hyperref}
\documentclass[12pt]{beamer}

\usetheme{TUDo}
\usefonttheme{professionalfonts}

\usepackage{fontspec}
\usepackage{polyglossia}
\setmainlanguage{german}

\usepackage{graphicx}
\usepackage{tikz}

\title[Toolbox 2015]{Toolbox-Workshop 2015}
\subtitle{Python, Unix-Shell, Make, Git, \LaTeX{}}
\date{5.–16.~Oktober 2015}
\institute[Pep et al.\ e.V.]{\includegraphics[height=0.5cm]{logos/pep.pdf}}
\titlegraphic{../python/muon_plot.png}
\author{PeP et al. e.\,V.}

\begin{document}

\maketitle

\begin{frame}
  \begin{center}
  \Large Praktikums-Versuche richtig auswerten!\\
  5. – 9.~Oktober 2015, 13–17~Uhr
  \end{center}
  \begin{tikzpicture}[remember picture, overlay]
    \tikzset{shift={(current page.center)}}
    \node[anchor=center] at (-3,2) {%
      \includegraphics[height=1cm]{logos/python.pdf}%
    };
    \node[anchor=center] at (3,2) {%
      \includegraphics[height=1.25cm]{logos/git.png}%
    };
    \node[anchor=center] at (4,-2) {%
      \includegraphics[height=1cm]{logos/numpy.png}%
    };
    \node[anchor=center] at (-2.5,-2) {%
      \includegraphics[height=1cm]{logos/matplotlib.png}%
    };
    \node[anchor=center] at (0,-3.5) {%
      \includegraphics[height=1.25cm]{logos/uncertainties.png}%
    };
  \end{tikzpicture}
\end{frame}

\begin{frame}
  \begin{center}
    Wissenschaftliche Arbeiten verfassen mit \\[0.5\baselineskip]
    \textrm{\Huge\LaTeX{}}\\[0.5\baselineskip]
    12. – 16.~Oktober 2015, 13–17~Uhr
  \end{center}
\end{frame}

\begin{frame}
  \begin{center}
    \Large
    Software installieren\\[0.5\baselineskip]
    → \textcolor{blue!70!black}{\url{http://toolbox.pep-dortmund.org/install.html}}\\[0.5\baselineskip]
    Laptops mitbringen / uns ansprechen\\[0.5\baselineskip]
    Umfrage ausfüllen
  \end{center}
\end{frame}
\begin{frame}
  \Huge\centering
  \textcolor{red!70!black}{Fragen?}
\end{frame}
\end{document}
